\section{The set of degree d surfaces including k-crosses}

It is natural to ask about the set of degree $d$ surfaces including $k$-crosses. Let $Z_k$ be a k-cross lying in a plane $H$, which we will assume is the plane $\mathbb{P}^2_{x,y}$ by an automorphism of $\mathbb{P}^3$. 

\begin{proposition}
    Let $A_k = \{S_d \subseteq \mathbb{P}^3 \mid Z_k \subset S_d\} \subseteq \{S_d \subseteq \mathbb{P}^3\}$ be the set of degree $d$ surfaces containing $Z_k$. Then $\dim{A_k} = \binom{d+1}{2} d-k+1$.
\end{proposition}

\begin{proof}
    $S_d = V(f)$ for some $f \in \mathbb{C}[x_0, x_1, x_2, x_3]_d$. Furthermore, since $Z_k$ is included in the $x,y$ plane, we write $f$ as $f = g(x,y) + zh(x,y,z)$, seperating the part of $f$ that vanishes on the plane. Since $Z_k \subseteq S_d$, we have that $g(x,y)$ vanishes on $Z_k$. That $C_k \subseteq S_d$ is the same as saying that $f\mid_{C_k} = 0$, which is the same as saying that $g(x,y) \mid_{C_k} = 0$. This is because $zh(x,y,z) \mid_{C_k} = 0$ since $C_k \subseteq \mathbb{P}^2_{x,y} = V(z)$. This is implies that $l_1l_2 \cdots l_k \mid g(x,y)$, meaning that $g(x,y) = l_1l_2 \dots l_k \cdot j(x,y)$ for some $j(x,y) \in \mathbb{C}[x,y]_{d-k}$. Thus we can write $f$ as $f = l_1l_2 \cdots l_k \cdot j(x,y) + zh(x,y,z)$, where $h(x,y,z) \in \mathbb{C}[x,y,z]_{d-1}$.

    Finishing the argument, we get that $\dim{A_k} = \dim \langle j(x,y), h(x,y,z) \rangle = \binom{d-k+1}{1} + \binom{d+1}{2}$. 
\end{proof}

\section{Surfaces including k-gons}

In this section, we will consider $k$-gons in $\mathbb{P}^3$ and their relationship with hypersurfaces.

\begin{definition}
    A $k$-gon is a union of $k$ lines $l_1, l_2, \dots, l_k$ such that $l_i$ intersects $l_{i+1}$ for $1 \leq i < k$, and $l_k$ intersects $l_1$. We will write $C_k = \bigcup_{i=1}^k l_i$ and denote the intersection points as $p_i = l_i \cap l_{i+1}$ for $1 \leq i < k$, and $p_k = l_k \cap l_1$.
\end{definition}

\begin{proposition}
    Let $C_k = \bigcup_{i=1}^k l_i$ be a $k$-gon in $\mathbb{P}^3$. Then we can find a hypersurface of degree $d$ such that $C_k \subseteq S_d$ when $k \leq \frac{(d+3)(d+2)(d+1)}{6d}$.
\end{proposition}

\begin{proof}

    We can look at the zero cohomology of the ideal sheaf $\mathcal{I}_{C_k}$ of the $k$-gon to see if the $k$-gon can be contained in a hypersurface. We will denote the union of the lines as $C_k = \bigcup_{i=1}^k l_i$. 

    We have the short exact sequence
    \[0 \to \mathcal{I}_{C_k} \to \mathcal{O}_{\mathbb{P}^3}\to \mathcal{O}_{C_k} \to 0.\]

    We twist this by $\mathcal{O}_{\mathbb{P}^3}(d)$ to get
    \[0 \to \mathcal{I}_{C_k}(d) \to \mathcal{O}_{\mathbb{P}^3}(d) \to \mathcal{O}_{C_k}(d) \to 0.\]

    From this we get the long exact sequence in cohomology:

    \begin{tikzcd}
    0 \rar & H^0(\mathbb{P}^3, \mathcal{I}_{C_k}) \rar & H^0(\mathbb{P}^3, \mathcal{O}_{\mathbb{P}^3}) \rar &  H^0(C_k, \mathcal{O}_{C_k}) \ar[out=-20, in=160]{dll} \\
    & H^1(\mathbb{P}^3, \mathcal{I}_{C_k}) \rar & H^1(\mathbb{P}^3, \mathcal{O}_{\mathbb{P}^3}) \rar & H^1(C_k, \mathcal{O}_{C_k}) \ar[out=-20, in=160]{dll} \\
    & H^2(\mathbb{P}^3, \mathcal{I}_{C_k}) \rar & H^2(\mathbb{P}^3, \mathcal{O}_{\mathbb{P}^3}) \rar & H^2(C_k, \mathcal{O}_{C_k}) \ar[out=-20, in=160]{dll}\\
    & H^3(\mathbb{P}^3, \mathcal{I}_{C_k}) \rar & H^3(\mathbb{P}^3, \mathcal{O}_{\mathbb{P}^3}) \rar & H^3(C_k, \mathcal{O}_{C_k}) \rar & 0
    \end{tikzcd}

    We know that $H^i(\mathbb{P}^3, \mathcal{O}_{\mathbb{P}^3}) = 0$ for $i > 0$. This means we can focus on the first part of the sequence:

    \[0 \to H^0(\mathbb{P}^3, \mathcal{I}_{C_k}) \to H^0(\mathbb{P}^3, \mathcal{O}_{\mathbb{P}^3}) \to H^0(C_k, \mathcal{O}_{C_k}) \to H^1(\mathbb{P}^3, \mathcal{I}_{C_k}) \to 0.\] 


    Since $h^0$ is additive, we get

    \[
    h^0(\mathbb{P}^3, \mathcal{I}_{C_k}) = h^0(\mathbb{P}^3, \mathcal{O}_{\mathbb{P}^3}) + h^1(\mathbb{P}^3, \mathcal{I}_{C_k}) - h^0(C_k, \mathcal{O}_{C_k}). 
    \]

    By Thm 18.4 in (EO) we get that $h^0(\mathbb{P}^3, \mathcal{O}_{\mathbb{P}^3}) = \binom{d+3}{3}$. $h^0(C_k, \mathcal{O}_{C_k}(d))$ can be computed by considering the short exact sequence: 
    \[
    0 \to \mathcal{O}_{C_k}(d) \to \bigoplus_{i=1}^k \mathcal{O}_{l_i}(d) \to \bigoplus_{i=1}^n \mathcal{O}_{p_i} \to 0.
    \]

    This gives us $h^0(C_k, \mathcal{O}_{C_k}) = k \cdot \binom{d+1}{1} - k = kd$.

    \[
    h^0(\mathbb{P}^3, \mathcal{I}_{C_k}) \geq \binom{d+3}{3} - kd.
    \]

    Meaning we can guarantee a degree $d$ surface when $\binom{d+3}{3} - kd \geq 0$, or equivalently when $k \leq \frac{(d+3)(d+2)(d+1)}{6d}$.

\end{proof}

To get an exact value for $h^0(\mathbb{P}^3, \mathcal{I}_{C_k})$, we would need to compute $h^1(\mathbb{P}^3, \mathcal{I}_{C_k})$. This issue will be explored in section ???. 


\begin{remark}
    Note that this argument does not guarantee that the surface is smooth, only that it exists.
\end{remark}

\section{Computational results}

The results above are supported by computational experiments in Macaulay2, where we have tested the existence of smooth hypersurfaces including $k$-gons for various degrees $d$ and number of lines $k$. The code used for these experiments is included in the following snippet. We will discuss the details of the code below, along with its results.

\lstinputlisting[caption={k-gon}, label={lst:k-gon}]{../Progging/ngon.m2}

We begin by defining a function to pick random points in our quotient ring, ensuring they are not zero. We then gather the points into pairs, so that we can later define lines between these. We then define a function for turning to points into a matrix, so that we can define the ideal by taking its minors. The function will turn two points $(a,b,c,d), (d,e,f,g)$ into a matrix $\begin{pmatrix}
    x_0 & x_1 & x_2 & x_3 \\
    a & b & c & d \\
    d & e & f & g 
\end{pmatrix}$, meaning we can define the ideal of the line between the two points by taking the $3 \times 3$ minors of this matrix. Next we take the intersection of all the line ideals, to get the ideal of $Z = \bigcup l_i$, and we the choose a random member of $I_Z(d)$. We then compute if this defines an empty, singular, or smooth surface. \\

\section{k-chains and special cases}

\begin{definition}
    A k-chain is a union of k lines $l_1, l_2, \dots, l_k$ such that $l_i$ intersects $l_{i+1}$ for $1 \leq i < k$. We will write $C_k = \bigcup_{i=1}^k l_i$ and denote the intersection points as $p_i = l_i \cap l_{i+1}$ for $1 \leq i < k$.
\end{definition}

Notice that a k-chain is a k-gon without the line $l_k$ intersecting $l_1$. One can then modify the code above to create k-chains instead of k-gons, by excluding the final pair of points. The results from this are similar to those of k-gons, in that for higher degrees, we find smooth surfaces, up until we find no surfaces. There are however interesting exceptions for degree 2 and 3. Here we get that for 4 lines and 6 lines, we get singular surfaces for 4 and 7 lines, respectively. We will handle the cases seperately.

\begin{proposition}
    For a general 4-chain in $\mathbb{P}^3$, we cannot find a smooth quadric surface $ S_2 \subset \mathbb{P}^3$ such that it includes the 4-chain.
\end{proposition}

\begin{proof}
    First, assume that $S_2$ is smooth. Since $S_2$ is a smooth quadric surface in $\mathbb{P}^3$, it is of the form $S_2 = \mathbb{P}^1 \times \mathbb{P}^1$. This means that $S_2$ has two rulings, meaning that there are two families of lines on $S_2$, such that any two lines in the same family do not intersect, and any two lines in different families intersect at exactly one point. Denote the two families as $F_1$ and $F_2$. Then we can see that a 4-chain must consist of two lines from $F_1$ and two lines from $F_2$, since any two lines in the same family do not intersect. This means that the 4-chain must be of the form $l_1, l_2, l_3, l_4$ where $l_1, l_3 \in F_1$ and $l_2, l_4 \in F_2$. However, this means that $l_4$ intersects $l_1$, since they are in different families, meaning that the 4-chain is in fact a 4-gon. This is a contradiction, since we assumed that we had a general 4-chain. Thus we cannot find a smooth quadric surface $ S_2 \subset \mathbb{P}^3$ such that it includes the 4-chain.
\end{proof}

\begin{proposition}
    For a general 4-chain in $\mathbb{P}^3$, we can find a singular quadric surface $ S_2 \subset \mathbb{P}^3$ such that it includes the 4-chain.
\end{proposition}
\begin{proof}
    Find a nice example of a singular quadric surface including a 4-chain.
\end{proof}

Now we will look at the the case of 7-chains in cubic surfaces. 

\begin{proposition}
    For a general 7-chain in $\mathbb{P}^3$, we cannot find a smooth cubic surface $ S_3 \subset \mathbb{P}^3$ such that it includes the 7-chain.
\end{proposition}

\begin{proof}
    First, the dimension of the space of cubic surfaces in $\mathbb{P}^3$ is given by the number of monomials of degree 3 in 4 variables, which is $\binom{3+4-1}{4-1} - 1 = \binom{6}{3} - 1 = 19$. Thus, the space of cubic surfaces in $\mathbb{P}^3$ has dimension 19. 

    Now, the linear system of 7-chains in $\mathbb{P}^3$ has dimension decided by the number of conditions imposed by the 7 lines. Recall that the set of lines in $\mathbb{P}^3$ is parametrized by the Grassmannian variety $\mathrm{Gr}(2,4)$, which is a 4-dimensional variety. Therefore the first line imposes 4 conditions, the second line imposes 3 conditions (since it must intersect the first line), the third line imposes 3 conditions (since it must intersect the second line), and so on. Thus, the total number of conditions imposed by the 7 lines is $4 + 3 + 3 + 3 + 3 + 3 + 3 = 22$. This being of a higher dimension than the space of cubic surfaces, we do not expect to find a cubic surface including our 7-chain.
\end{proof}


Note that if we instead asked whether a smooth cubic surface in $\mathbb{P}^3$ can include a 7-chain, the situation is different. In this case, we can indeed find a 7-chain such that it is included in the surface. In fact, we know we could find a 12-chain in this case, and a 13-gon. We gather show this in the following proposition. 

\begin{proposition}
    For a smooth cubic surface in $\mathbb{P}^3$, we can find a 13-gon such that it is included in the surface.
\end{proposition}

\begin{proof}
    A smooth cubic surface in $\mathbb{P}^3$ is isomorphic to the blow-up of $\mathbb{P}^2$ in 6 points. The surface then includes 27 lines, which are precisely the exceptional divisors of the blow-up, the strict transforms of the 15 lines through pairs of points, and the strict transforms of the 6 conics through 5 of the points. We can then find a 13-gon by considering the following lines: Let $E_1, E_2, E_3, E_4, E_5, E_6$ be the exceptional divisors, let $L_{ij}$ be the strict transform of the line through points $p_i$ and $p_j$, and let $C_i$ be the strict transform of the conic through all points except $p_i$. Then we can consider the 13-gon given by the lines $E_1, L_{12}, E_2, L_{23}, E_3, L_{34}, E_4, L_{45}, E_5, L_{56}, E_6, C_1, L_{16}$. This is indeed a 13-gon since each line intersects exactly two other lines in the set.
\end{proof}

By removing one of the lines from the 13-gon, we immediately get the following corollary:
\begin{corollary}
    For a smooth cubic surface in $\mathbb{P}^3$, we can find a 12-chain such that it is included in the surface.
\end{corollary}